\section*{Capítulo \ref{ch:g12:diffeq} --- Equações diferenciais}

\begin{solution}{exo:g12:diffeq:1}
\emph{(a)} $y = C\eu^{3x}$; $y(0) = 2$ dá $y = 2\eu^{3x}$.

\emph{(b)} $y' = -\frac12 y$, logo $y = C\eu^{-x/2}$; $y(0) = -1$ dá
$y = -\eu^{-x/2}$.

\emph{(c)} Equilíbrio $y_p = 5$; $y = C\eu^{-x} + 5$; $y(0) = 0$ dá
$C = -5$: $y = 5\left(1 - \eu^{-x}\right)$.
\end{solution}

\begin{solution}{exo:g12:diffeq:2}
$N' = kN$, logo $N(t) = N_0 \eu^{kt}$. Dobrar em $3$ horas significa
$\eu^{3k} = 2$, \emph{isto é},
\[
k = \frac{\ln 2}{3} \approx 0.231\ \text{h}^{-1}.
\]
\end{solution}

\begin{solution}{exo:g12:diffeq:3}
$y_p'(x) = \eu^x + x\eu^x = y_p(x) + \eu^x$: $y_p$ é uma solução particular
(este é o caso ressonante do \cref{met:g12:diffeq:particular}). Pela
\cref{prop:g12:diffeq:general}, as soluções são
$y = C\eu^{x} + x\eu^{x} = (C + x)\,\eu^x$, $C \in \R$.
\end{solution}

\begin{solution}{exo:g12:diffeq:4}
Tente $y_p = \beta\eu^x$: $\beta\eu^x = -2\beta\eu^x + \eu^x$ dá
$3\beta = 1$, logo $y_p = \frac13\eu^x$. Solução geral
$y = C\eu^{-2x} + \frac13\eu^{x}$; a condição $y(0) = 1$ dá
$C = \frac23$:
\[
y(x) = \frac{2}{3}\,\eu^{-2x} + \frac{1}{3}\,\eu^{x}.
\]
\end{solution}

\begin{solution}{exo:g12:diffeq:5}
$N(t) = N_0\,\eu^{-\lambda t}$ com $\lambda = \frac{\ln 2}{5730}$
(\cref{ex:g12:exp:halflife}). Resolvemos $\eu^{-\lambda t} = 0.6$:
\[
t = \frac{\ln(1/0.6)}{\lambda} = 5730\,\frac{\ln(5/3)}{\ln 2}
\approx 5730 \times \frac{0.5108}{0.6931} \approx 4220 \text{ anos}.
\]
\end{solution}

\begin{solution}{exo:g12:diffeq:6}
\emph{1.} O sal entra a $0.2 \times 5 = 1$ kg/min. A saída leva
concentração $\frac{m}{100}$ kg/L a $5$ L/min, \emph{isto é},
$\frac{m}{20}$ kg/min. Logo $m' = 1 - \frac{m}{20}$.

\emph{2.} Equilíbrio $m_p = 20$; $m(t) = 20 + C\eu^{-t/20}$, e $m(0) = 0$
dá $C = -20$:
\[
m(t) = 20\left(1 - \eu^{-t/20}\right) \xrightarrow[t\to+\infty]{} 20 \text{ kg}.
\]
A longo prazo, a concentração do tanque iguala a da salmoura que entra:
$0.2$ kg/L $\times$ $100$ L $= 20$ kg.
\end{solution}

\begin{solution}{exo:g12:diffeq:7}
\emph{1.} $v' = g - \frac{k}{m} v$ com
$\frac km = \frac{16}{80} = 0.2$. Equilíbrio
$v_\infty = \frac{mg}{k} = \frac{9.8}{0.2} = 49$ m/s;
$v(t) = 49 + C\eu^{-0.2t}$, e $v(0) = 0$ dá
\[
v(t) = 49\left(1 - \eu^{-0.2 t}\right).
\]

\emph{2.} Velocidade terminal $49$ m/s ($\approx 176$ km/h). Queremos
$1 - \eu^{-0.2t} = 0.95$, \emph{isto é}, $\eu^{-0.2t} = 0.05$:
\[
t = \frac{\ln 20}{0.2} \approx \frac{3.00}{0.2} \approx 15 \text{ s}.
\]
\end{solution}

\begin{solution}{exo:g12:diffeq:8}
\emph{1.} $z = \frac1y$ dá
$z' = -\frac{y'}{y^2} = -\frac{y(1-y)}{y^2} = -\frac{1-y}{y}
= -\frac1y + 1 = -z + 1$.

\emph{2.} Equilíbrio $z_p = 1$, logo $z(t) = 1 + C\eu^{-t}$. De
$y(0) = \frac{1}{10}$ vem $z(0) = 10$, logo $C = 9$ e
\[
y(t) = \frac{1}{1 + 9\,\eu^{-t}} .
\]

\emph{3.} Quando $t \to +\infty$, $9\eu^{-t} \to 0$ e $y(t) \to 1$. A
curva é a clássica curva logística em S (\emph{sigmoide}): crescimento
lento no início ($y$ pequeno, $y' \approx y$), crescimento \omterm{def:g10:functions:extrema}{máximo} quando
$y = \frac12$ (onde $y' = y(1-y)$ é \omterm{def:g10:functions:extrema}{máximo}) e depois saturação rumo à
capacidade de suporte $1$.
\end{solution}

\begin{solution}{exo:g12:diffeq:9}
Pelo \cref{thm:g12:diffeq:affine}, $y(x) = C\eu^{ax} - \frac ba$. Logo
\[
u_{n+1} = C\eu^{a(n+1)} - \frac ba
= \eu^{a}\left(C\eu^{an} - \frac ba\right) + \frac ba\left(\eu^a - 1\right)
= q\,u_n + r
\]
com $q = \eu^{a}$ e $r = \frac{b}{a}\left(\eu^{a} - 1\right)$. A condição
$\abs q < 1$ significa $\eu^a < 1$, \emph{isto é}, $a < 0$: exatamente a
condição sob a qual as soluções da \omterm{def:g12:diffeq:ode}{equação diferencial} convergem para o
equilíbrio $-\frac ba$ --- e, de fato, o \omterm{pb:g10:functions:1}{ponto fixo} da recorrência é
$\frac{r}{1 - q} = -\frac{b}{a}$.
\end{solution}

\begin{solution}{pb:g12:diffeq:1}
\textbf{1.} $y = 2\eu^{3t}$. Para a segunda: equilíbrio $3$,
logo $y = 3 + C\eu^{-2t}$ com $y(0) = 0$: $C = -3$:
$y = 3\left(1 - \eu^{-2t}\right)$.

\textbf{2.} $\left(y\eu^{-at}\right)' = y'\eu^{-at} -
ay\eu^{-at} = (y' - ay)\eu^{-at} = 0$: o produto é uma
constante $C$, logo $y = C\eu^{at}$ --- nenhuma solução
escapa.

\textbf{3.} Equilíbrio: $y \equiv 3$. Toda outra solução é
$3 + C\eu^{-2t}$ com $C \neq 0$: a exponencial morre e a
solução desliza até $3$ --- um \omterm{pb:g10:functions:1}{ponto fixo} estável, o gêmeo
\omterm{def:g12:limcont:continuity}{contínuo} dos \omterm{pb:g10:functions:1}{pontos fixos} das recorrências do
\cref{pb:g11:seq:1}.

\textbf{4.} $y = y_0\eu^{-t/\tau}$ cai à metade quando
$\eu^{-t/\tau} = \frac12$: $t = \tau\ln 2$.

\textbf{5.} Todas as soluções são translações verticais do
decaimento rumo a $3$: as que começam acima caem até $3$, as
que começam abaixo sobem até $3$, e a solução constante $3$
fica parada --- um funil em torno do equilíbrio.

\textbf{6.} $\lambda = \frac{\ln 2}{5730} \approx 1.21 \times
10^{-4}$ por ano.

\textbf{7.} $\eu^{-\lambda t} = 0.2$:
$t = \frac{\ln 5}{\lambda} \approx 13\,300$ anos.

\textbf{8.} $t = \frac{\ln(1/0.15)}{\lambda} \approx 15\,700$
anos: os touros de Lascaux são do fim da era glacial.

\textbf{9.} $t = \frac{\ln(1/0.53)}{\lambda} \approx 5\,250$
anos: dentro de uma vida humana do valor dos arqueólogos --- o
relógio funciona.

\textbf{10.} Dez meias-vidas deixam
$2^{-10} \approx 0.1\,\%$: no limite da medição. Após $65$
milhões de anos a fração é
$2^{-65\,000\,000/5\,730} \approx 2^{-11\,344}$: nenhum átomo
do estoque original resta em qualquer fóssil --- os dinossauros
são datados por relógios mais lentos (potássio-argônio e
parentes).

\textbf{11.} $52\,\%$ dá $\approx 5\,405$ anos e $54\,\%$ dá
$\approx 5\,095$: o $\pm 1\,\%$ da química vira cerca de
$\pm 155$ anos de história --- precisão no laboratório é
precisão na etiqueta do museu.

\textbf{12.} $z = T - T_a$ satisfaz $z' = -kz$:
$z = (T_0 - T_a)\eu^{-kt}$, e $T = T_a + z$.

\textbf{13.} A diferença inicial é $T_0 - T_a = 70$; após cinco
minutos ela é $70 - 20 = 50$, logo $50 = 70\eu^{-5k}$:
$k = \frac{\ln(70/50)}{5} \approx 0.067$ por minuto. Para
beber: $55 - 20 = 35 = 70\eu^{-kt}$:
$t = \frac{\ln 2}{k} \approx 10.3$ minutos.

\textbf{14.} A taxa de perda é proporcional à diferença
$T - T_a$: café fervendo sangra calor mais depressa.
Acrescentar o leite \emph{agora} corta a diferença
imediatamente, de modo que se perde menos calor ao longo dos
dez minutos; deixar para o fim permite que o café esfrie a toda
velocidade antes. Para a xícara mais quente: leite primeiro.
(Anfitriões impacientes têm a termodinâmica ao contrário.)

\textbf{15.} Diferenças em relação ao ambiente: à meia-noite
$10$, uma hora depois $8$: $\eu^{-k} = 0.8$:
$k = \ln\frac{10}{8} \approx 0.223$ por hora. Na morte a
diferença era $17$; ela decaiu até $10$ à meia-noite:
$s = \frac{\ln(17/10)}{k} \approx 2.4$ horas. Hora da morte:
por volta das $21{:}40$.

\textbf{16.} Uma sala aquecida ou com corrente de ar ($T_a$ não
constante) entorta o relógio para qualquer lado; roupas ou a
massa corporal mudam $k$ (as tabelas do perito corrigem isso)
--- um $k$ errado reescala todo o \omterm{def:g10:numbers:interval}{intervalo}; e um corpo movido
de outro lugar reinicia $T_a$ no meio do decaimento, forjando
uma morte mais cedo ou mais tarde. A \omterm{def:g10:algebra:equation}{equação} é honesta; o risco
está nas hipóteses.

\textbf{17.} $z = \frac1y$ obedece a $z' = 1 - z$:
$z = 1 + 9\eu^{-t}$ (de $z(0) = 10$), logo
$y = \frac{1}{1 + 9\eu^{-t}}$. Inflexão em $y = \frac12$:
$9\eu^{-t} = 1$: $t = \ln 9 \approx 2.2$ --- a data de virada
da epidemia, direto da fórmula.

\textbf{18.} Velocidade terminal: $v' = 0$ em $v = 20$ m/s.
Solução: $v(t) = 20\left(1 - \eu^{-t/2}\right)$. Então
$0.95$: $\eu^{-t/2} = 0.05$: $t = 2\ln 20 \approx 6$ s --- as
gotas de chuva alcançam sua velocidade de cruzeiro em segundos,
e é por isso que a chuva não mata.

\textbf{19.} Regime permanente: $c = 8$. Metade disso:
$4 = 8\left(1 - \eu^{-t/2}\right)$: $\eu^{-t/2} = \frac12$:
$t = 2\ln 2 \approx 1.4$ hora. O soro é o gêmeo \omterm{def:g12:limcont:continuity}{contínuo} do
empréstimo: entrada constante, saída proporcional --- o
\cref{exo:g12:diffeq:9} torna o dicionário exato.

\textbf{20.} Decaimento ($a < 0$, $b = 0$), resfriamento
($y = T - T_a$), queda ($a < 0$, $b > 0$: \omterm{def:g10:numbers:approx}{aproximação} do
equilíbrio por baixo), dosagem (a mesma coisa, numa veia): uma
\omterm{def:g10:algebra:equation}{equação}, quatro mundos. O ciclo: enunciar a lei da mudança;
escrever a \omterm{def:g10:algebra:equation}{equação}; resolver
(\cref{met:g12:diffeq:solve}); calibrar as constantes com
medições; prever; e então auditar as hipóteses (questão 16). A
oscilação exige $y'' = -\omega^2 y$ --- encontrada, resolvida e
posta a balançar no \cref{pb:g12:trigo:1}.
\end{solution}
